\renewcommand{\Anno}{2007}

\renewcommand{\NumeroProblema}{2}

\renewcommand{\Testo}{Sono date $k$ scatole ($k\geq 3$), numerate
da 1 a $k$.  Tre di queste scatole, con i numeri \emph{consecutivi},
pesano 99 grammi, le restanti pesano 100 grammi.  Una \emph{pesata}
consiste nello scegliere alcune scatole a piacere e porle su una
bilancia, che fornisce il peso totale.
	\begin{enumerate}
		\item Nel caso $k=6$, mostrare che \`{e} possibile individuare
		le 3 scatole da 99 grammi mediante una sola pesata.
	
		\item  Nel caso $k=12$, mostrare che \`{e} possibile individuare
		le 3 scatole da 99 grammi mediante 2 pesate.
	
		\item Determinare se nel caso $k=18$ \`{e} ancora possibile
		individuare le 3 scatole da 99 grammi mediante 2 pesate.
	\end{enumerate}
}

\renewcommand{\Risposta}{
\begin{enumerate}
	\item  \ldots
		
	\item  \ldots
		
	\item \`{E} possibile
\end{enumerate}
}

\renewcommand{\Aiutino}{
\begin{enumerate}
	\renewcommand{\labelenumi}{(\alph{enumi})} 
	\item  Pesiamo le prime 3 scatole \ldots

	\item  Pesiamo le prime 6 scatole \ldots
	
	\item Al primo colpo pesare le scatole 3, 4, 5, 6, 13, 14, 15, 16.
\end{enumerate}
}

\renewcommand{\SoluzioneUfficiale}{\TitoloSoluzione\DomandaA Pesiamo
le scatole 1, 2, 3.  Si presentano 4 casi, a seconda che tra
queste ve ne siano 0, 1, 2 o 3 da 99 grammi.
\begin{center}
	\psset{unit=3.5ex}
	\pspicture(1,0)(7,1)
	\psframe[linestyle=none,fillstyle=solid,fillcolor=green](1,0)(4,1)
	\multido{\n=1+1}{6}{
	\psframe(\n,0)(!\n\space 1 add 1)
	\rput(!\n\space 0.5 add 0.5){$\n$}
	}
	\endpspicture
\end{center}

\begin{itemize}
	\item  Se tra le prime 3 non ci sono scatole da 99 grammi, allora 
	le scatole da 99 grammi sono in posizione 4--5--6.

	\item  Se tra le prime 3 c'\`{e} una scatola da 99 grammi, allora 
	le scatole da 99 grammi sono in posizione 3--4--5.

	\item  Se tra le prime 3 ci sono 2 scatole da 99 grammi, allora 
	le scatole da 99 grammi sono in posizione 2--3--4.

	\item  Se tutte le prime 3 ci sono 3 scatole da 99 grammi, allora 
	ovviamente le scatole da 99 grammi sono in posizione 1--2--3.
\end{itemize}

\DomandaB Al primo colpo pesiamo insieme le prime 6 scatole.
\begin{center}
	\psset{unit=3.5ex}
	\pspicture(1,0)(13,1)
	\psframe[linestyle=none,fillstyle=solid,fillcolor=green](1,0)(7,1)
	\multido{\n=1+1}{12}{
	\psframe(\n,0)(!\n\space 1 add 1)
	\rput(!\n\space 0.5 add 0.5){$\n$}
	}
	\endpspicture
\end{center}
\begin{itemize}
	\item Se tra le prime 6 scatole non ve ne sono da 99 grammi, vuol
	dire che le scatole da 99 grammi si trovano nelle ultime 6
	posizioni, e quindi si possono individuare con una sola pesata
	come indicato nella prima parte del problema.

	\item Se tra le prime 6 scatole pesate ce n'\`{e} una da 99 grammi, vuol
	dire che le scatole da 99 grammi sono in posizione 6--7--8.

	\item Se tra le prime 6 scatole pesate ce ne sono 2 da 99 grammi, vuol
	dire che le scatole da 99 grammi sono in posizione 5--6--7.

	\item Se tra le prime 6 scatole pesate ce ne sono 3 da 99 grammi,
	vuol dire che le scatole da 99 grammi sono tutte nelle prime 6
	posizioni, dunque si possono individuare con la restante pesata
	come indicato nella prima parte del problema.
\end{itemize}

\DomandaC Pesiamo insieme 8 scatole, e cio\`{e} quelle in posizione 3,
4, 5, 6, 13, 14, 15, 16.
\begin{center}
	\psset{unit=3.5ex}
	\pspicture(1,0)(19,1)
	\psframe[linestyle=none,fillstyle=solid,fillcolor=green](3,0)(7,1)
	\psframe[linestyle=none,fillstyle=solid,fillcolor=green](13,0)(17,1)
	\multido{\n=1+1}{18}{
	\psframe(\n,0)(!\n\space 1 add 1)
	\rput(!\n\space 0.5 add 0.5){$\n$}
	}
	\endpspicture
\end{center}

\begin{itemize}
	\item Se tra le 8 scatole pesate non ve ne sono da 99 grammi, vuol
	dire che le scatole da 99 grammi si trovano nelle 6 posizioni che
	vanno dalla 7 alla 12, e quindi si possono individuare con una
	sola pesata come indicato nella prima parte del problema.

	\item Se tra le 8 scatole pesate ce n'\`{e} una da 99 grammi, vuol
	dire che le scatole da 99 grammi possono essere in 4 posizioni:
	1--2--3, 6--7--8, 11--12--13, 16--17--18.
\begin{center}
	\psset{unit=3.5ex}
	\pspicture(1,0)(19,1)
	\psframe[linestyle=none,fillstyle=solid,fillcolor=yellow](1,0)(4,1)
	\psframe[linestyle=none,fillstyle=solid,fillcolor=yellow](6,0)(9,1)
	\psframe[linestyle=none,fillstyle=solid,fillcolor=yellow](11,0)(14,1)
	\psframe[linestyle=none,fillstyle=solid,fillcolor=yellow](16,0)(19,1)
	\multido{\n=1+1}{18}{
	\psframe(\n,0)(!\n\space 1 add 1)
	\rput(!\n\space 0.5 add 0.5){$\n$}
	}
	\endpspicture
\end{center}
	
	Per decidere tra queste pesiamo insieme
	6, 11, 12, 16, 17, 18 (o comunque 0 scatole prese dal primo
	gruppo, 1 dal secondo, 2 dal terzo, 3 dal quarto).  A seconda che
	in questa seconda pesata troviamo 0, 1, 2, 3 scatole da 99 grammi,
	il gruppo buono sar\`{a}, rispettivamente, il primo, secondo,
	terzo, quarto.

	\item Se tra le 8 scatole pesate ce ne sono 2 da 99 grammi, vuol
	dire che le scatole da 99 grammi possono essere in 4 posizioni:
	2--3--4, 5--6--7, 12--13--14, 15--16--17.  Come nel caso
	precedente pesiamo insieme 0 scatole prese dal primo gruppo, 1 dal
	secondo, 2 dal terzo, 3 dal quarto (ad esempio 5, 12, 13, 15, 16,
	17) e concludiamo allo stesso modo.

	\item Se tra le 8 scatole pesate ce ne sono 2 da 99 grammi, vuol
	dire che queste possono essere in 4 posizioni: 3--4--5, 4--5--6,
	13--14--15, 14--15--16.  Pesiamo ora insieme 6, 14, 15, 16.  \`{E}
	facile verificare che, a seconda che in questa pesata troviamo 0,
	1, 2, 3 scatole da 99, il gruppo buono sar\`{a}, rispettivamente,
	il primo, secondo, terzo, quarto.
\end{itemize}
}

\renewcommand{\SoluzioniAlternative}{\TitoloSoluzioneAlternativaA 

Una soluzione alternativa per la seconda domanda \`{e} la seguente. 
Al primo colpo pesiamo insieme le prime 3 e le ultime 3 scatole.
\begin{center}
	\psset{unit=3.5ex}
	\pspicture(1,0)(13,1)
	\psframe[linestyle=none,fillstyle=solid,fillcolor=green](1,0)(4,1)
	\psframe[linestyle=none,fillstyle=solid,fillcolor=green](10,0)(13,1)
	\multido{\n=1+1}{12}{
	\psframe(\n,0)(!\n\space 1 add 1)
	\rput(!\n\space 0.5 add 0.5){$\n$}
	}
	\endpspicture
\end{center}
\begin{itemize}
	\item Se tra le 6 scatole pesate non ve ne sono da 99 grammi, vuol
	dire che le scatole da 99 grammi si trovano nelle 6 posizioni che
	vanno dalla 4 alla 9, e quindi si possono individuare con una sola
	pesata come indicato nella prima parte del problema.

	\item Se tra le 6 scatole pesate ce n'\`{e} una da 99 grammi, vuol
	dire che le scatole da 99 grammi sono in posizione 3--4--5 oppure
	8--9--10.  A questo punto basta pesare, per esempio, la 3 per
	decidere tra le 2 alternative.

	\item Se tra le 6 scatole pesate ce ne sono 2 da 99 grammi, vuol
	dire che le scatole da 99 grammi sono in posizione 2--3--4 oppure
	9--10--11.  A questo punto basta pesare, per esempio, la 3 per
	decidere tra le 2 alternative.

	\item Se tra le 6 scatole pesate ce ne sono 3 da 99 grammi, vuol
	dire che le scatole da 99 grammi sono le prime 3 o le ultime 3.
	Ancora una volta basta pesare, per esempio, la 3 per decidere tra
	le 2 alternative.
\end{itemize}

Segnaliamo ora una soluzione alternativa alla terza domanda, in cui la
seconda pesata viene decisa a priori, e non dopo aver visto l'esito
della prima. Le 2 pesate sono le seguenti.
\begin{center}
	\psset{unit=3.5ex}
	\pspicture(1,0)(19,1)
	\psframe[linestyle=none,fillstyle=solid,fillcolor=green](3,0)(7,1)
	\psframe[linestyle=none,fillstyle=solid,fillcolor=green](13,0)(17,1)
	\multido{\n=1+1}{18}{
	\psframe(\n,0)(!\n\space 1 add 1)
	\rput(!\n\space 0.5 add 0.5){$\n$}
	}
	\endpspicture
\end{center}

\begin{center}
	\psset{unit=3.5ex}
	\pspicture(1,0)(19,1)
	\psframe[linestyle=none,fillstyle=solid,fillcolor=green](6,0)(7,1)
	\psframe[linestyle=none,fillstyle=solid,fillcolor=green](10,0)(13,1)
	\psframe[linestyle=none,fillstyle=solid,fillcolor=green](14,0)(19,1)
	\multido{\n=1+1}{18}{
	\psframe(\n,0)(!\n\space 1 add 1)
	\rput(!\n\space 0.5 add 0.5){$\n$}
	}
	\endpspicture
\end{center}

Nella seguente tabella le intestazioni delle colonne rappresentano i
possibili esiti della prima pesata (intesi come numero di scatole da
99 grammi presenti), le intestazioni delle righe rappresentano i
possibili esiti della seconda pesata, i contenuti delle caselle la
terna di scatole da 99 grammi nei vari casi.

\begin{center}
	\renewcommand{\arraystretch}{1.4}
	\begin{tabular}{|c||c|c|c|c|}
		\cline{2-5}
		\multicolumn{1}{c|}{} & 0 & 1 & 2 & 3  \\
		\hline
		\hline
		0 & 7--8--9 & 1--2--3 & 2--3--4 & 3--4--5  \\
		\hline
		1 & 8--9--10 & 6--7--8 & 5--6--7 & 4--5--6  \\
		\hline
		2 & 9--10--11 & 11--12--13 & 12--13--14 & 13--14--15  \\
		\hline
		3 & 10--11--12 & 16--17--18 & 15--16--17 & 14--15--16  \\
		\hline
	\end{tabular}
\end{center}

\`{E} chiaro che queste 2 pesate, ristrette alle prime 12 scatole,
forniscono pure una risposta alla seconda domanda del problema. 

}

\renewcommand{\Osservazioni}{\TitoloOsservazione Nella terza domanda
del problema le possibili configurazioni sono 16 (la prima scatola da
99 grammi pu\`{o} infatti trovarsi in una qualunque posizione dalla 1
alla 16).  Ogni pesata ha 4 esiti possibili (a seconda che si trovino
0, 1, 2, 3 scatole da 99 grammi).  Pertanto 2 pesate possono dare in
totale $4\cdot 4=16$ esiti distinti.  Questo mostra che il caso
presentato nella terza parte del problema \`{e} il caso limite con 2
pesate.  Se le scatole fossero 19 o pi\`{u}, le configurazioni
possibili sarebbero 17 o pi\`{u}, dunque non si potrebbero
discriminare con sole 2 pesate.

Analogamente 6 scatole (cio\`{e} 4 configurazioni possibili)
rappresentano il caso limite con una sola pesata.

Il problema si presta ad essere generalizzato al caso di $m$ pesate.
Il ragionamente appena presentato si estende facilmente mostrando che,
con un numero di scatole maggiore od uguale di $4^{m}+3$, non c'\`{e}
possibilit\`{a} di discriminare.  Quello che non sembra semplice \`{e}
costruire una soluzione generale con $m$ pesate nel caso di $4^{m}+2$
scatole (sempre che si possa fare\ldots).  }

\renewcommand{\Commenti}{
}